% Part: computability
% Chapter: recursive-functions
% Section: pr-functions

\documentclass[../../../include/open-logic-section]{subfiles}

\begin{document}

\olfileid{cmp}{rec}{prf}
\olsection{আদিম পুনরাবৃত্ত অপেক্ষক}

আদিম পুনরাবৃত্তি ও মিশ্রণ ব্যবহার করে বিদ্যমান অপেক্ষকগুলি থেকে কীভাবে
নতুন অপেক্ষক সংজ্ঞায়িত করা যায়, তা আবার লিখে রাখা যাক।

\begin{defn}
  \ollabel{defn:primitive-recursion} ধরা যাক, $f$ একটি $k$-স্থানীয়
  অপেক্ষক ($k\ge 1$) এবং $g$ একটি $(k+2)$-স্থানীয় অপেক্ষক। $f$ ও
  $g$ থেকে \emph{আদিম পুনরাবৃত্তির মাধ্যমে সংজ্ঞায়িত অপেক্ষক} হলো
  সমীকরণগুলি দিয়ে সংজ্ঞায়িত $(k+1)$-স্থানীয় অপেক্ষক~$h$:
  \begin{align*}
    h(x_0,\dots,x_{k-1},0) & =  f(x_0,\dots,x_{k-1}) \\
    h(x_0,\dots,x_{k-1},y+1) & = g(x_0,\dots,x_{k-1}, y, h(x_0,\dots,x_{k-1}, y))
  \end{align*}
\end{defn}

\begin{defn}
  \ollabel{defn:composition} ধরা যাক, $f$ একটি $k$-স্থানীয় অপেক্ষক,
  এবং $g_0$, \dots, $g_{k-1}$ হলো $k$টি অপেক্ষক, যাদের প্রত্যেকটি
  $n$-স্থানীয়। $f$ ও $g_0$, \dots,~$g_{k-1}$ থেকে \emph{মিশ্রণের
  মাধ্যমে সংজ্ঞায়িত অপেক্ষক} হলো নিম্নলিখিতভাবে সংজ্ঞায়িত
  $n$-স্থানীয় অপেক্ষক~$h$:
  \[
  h(x_0, \dots, x_{n-1}) =
  f(g_0(x_0, \dots, x_{n-1}), \dots, g_{k-1}(x_0, \dots, x_{n-1})).
  \]
\end{defn}

$\Succ$ এবং অভিক্ষেপ অপেক্ষকগুলি
\[
\Proj{n}{i}(x_0,\dots,x_{n-1}) = x_i,
\]
ছাড়াও, প্রত্যেক স্বাভাবিক সংখ্যা $n$ ও $i < n$-এর জন্য আমরা
অপেক্ষক~$\Zero(x) = 0$-কে আদিম পুনরাবৃত্ত অপেক্ষকগুলির অন্তর্ভুক্ত করব।

\begin{defn}
  আদিম পুনরাবৃত্ত অপেক্ষকসমূহের সেট হলো $\Nat^n$ থেকে $\Nat$-এ
  অপেক্ষকগুলির সেই সেট, যা নিচের শর্তগুলি দিয়ে আবেশনভিত্তিকভাবে
  সংজ্ঞায়িত:
  \begin{enumerate}
  \item $\Zero$ আদিম পুনরাবৃত্ত।
  \item $\Succ$ আদিম পুনরাবৃত্ত।
  \item প্রত্যেক অভিক্ষেপ অপেক্ষক $\Proj{n}{i}$ আদিম পুনরাবৃত্ত।
  \item যদি $f$ একটি $k$-স্থানীয় আদিম পুনরাবৃত্ত অপেক্ষক এবং $g_0$,
    \dots,~$g_{k-1}$ হলো $n$-স্থানীয় আদিম পুনরাবৃত্ত অপেক্ষক, তবে
    $g_0$, \dots,~$g_{k-1}$-এর সঙ্গে $f$-এর মিশ্রণ আদিম পুনরাবৃত্ত।
  \item যদি $f$ একটি $k$-স্থানীয় আদিম পুনরাবৃত্ত অপেক্ষক এবং $g$ একটি
    $k+2$-স্থানীয় আদিম পুনরাবৃত্ত অপেক্ষক হয়, তবে $f$ ও $g$ থেকে
    আদিম পুনরাবৃত্তির মাধ্যমে সংজ্ঞায়িত অপেক্ষকটি আদিম পুনরাবৃত্ত।
\end{enumerate}
\end{defn}

\begin{explain}
আরও সংক্ষেপে, আদিম পুনরাবৃত্ত অপেক্ষকসমূহের সেট হলো $\Zero$, $\Succ$
ও অভিক্ষেপ অপেক্ষকগুলি~$\Proj{n}{j}$-কে ধারণকারী এবং মিশ্রণ ও আদিম
পুনরাবৃত্তির অধীনে বদ্ধ ক্ষুদ্রতম সেট।

আদিম পুনরাবৃত্ত অপেক্ষকসমূহের সেটকে “পর্যায়”-এর সাহায্যেও বর্ণনা করা
যায়। $S_0$ দিয়ে প্রারম্ভিক অপেক্ষকগুলির সেট বোঝাই: $\Zero$, $\Succ$
ও অভিক্ষেপগুলি। এগুলিই পর্যায়~$0$-এর আদিম পুনরাবৃত্ত অপেক্ষক। কোনো
পর্যায় $S_i$ সংজ্ঞায়িত হয়ে গেলে, ইতিমধ্যে~$S_i$-তে থাকা অপেক্ষকগুলির
উপর একবার মিশ্রণ বা আদিম পুনরাবৃত্তি প্রয়োগ করে পাওয়া সব অপেক্ষকের
সেটকে $S_{i+1}$ বলা যাক। তাহলে
\[
S = \bigcup_{i \in \Nat} S_i
\]
হলো সব আদিম পুনরাবৃত্ত অপেক্ষকের সেট।
\end{explain}

$\Add$ যে একটি আদিম পুনরাবৃত্ত অপেক্ষক, তা যাচাই করা যাক।

\begin{prop}
  যোগের অপেক্ষক $\Add(x,y) = x+y$ আদিম পুনরাবৃত্ত।
\end{prop}

\begin{proof}
$\Add$-এর একটি আদিম পুনরাবৃত্ত সংজ্ঞা আমাদের ইতিমধ্যেই আছে, যা দুটি
অপেক্ষক $f$ ও~$g$-এর সাহায্যে দেওয়া এবং
\olref{defn:primitive-recursion}-এর বিন্যাসের সঙ্গে মেলে:
\begin{align*}
  \Add(x_0, 0) & = f(x_0) = x_0\\
  \Add(x_0, y+1) & = g(x_0, y, \Add(x_0, y)) =
  \Succ(\Add(x_0, y))
\end{align*}
অতএব $f$ ও $g$ আদিম পুনরাবৃত্ত হলেই $\Add$-ও আদিম পুনরাবৃত্ত।
$f(x_0) = x_0 = \Proj{1}{0}(x_0)$, এবং অভিক্ষেপ অপেক্ষকগুলি আদিম
পুনরাবৃত্ত হিসেবে গণ্য হয়, তাই $f$ আদিম পুনরাবৃত্ত। অপেক্ষক $g$ হলো
নিচের সমীকরণে সংজ্ঞায়িত তিন-স্থানীয় অপেক্ষক $g(x_0, y, z)$:
\[
g(x_0, y, z) = \Succ(z).
\]
এ থেকে এখনো বলা যায় না যে $g$ আদিম পুনরাবৃত্ত, কারণ $g$ ও $\Succ$
ঠিক একই অপেক্ষক নয়: $\Succ$ এক-স্থানীয়, আর $g$-কে তিন-স্থানীয় হতে
হবে। কিন্তু মিশ্রণের সাহায্যে আমরা $g$-কে “বিধিসম্মতভাবে” এভাবে
সংজ্ঞায়িত করতে পারি:
\[
g(x_0, y, z) = \Succ(\Proj{3}{2}(x_0, y, z))
\]
যেহেতু $\Succ$ ও $\Proj{3}{2}$ আদিম পুনরাবৃত্ত অপেক্ষক হিসেবে গণ্য হয়,
এবং অপেক্ষকটিকে আদিম পুনরাবৃত্ত অপেক্ষকগুলির মিশ্রণ দিয়ে সংজ্ঞায়িত করা যায়,
তাই $g$-ও আদিম পুনরাবৃত্ত।
\end{proof}

\begin{prop}
  \ollabel{prop:mult-pr}
  গুণের অপেক্ষক $\Mult(x,y) = x \cdot y$ আদিম পুনরাবৃত্ত।
\end{prop}

\begin{proof}
  অনুশীলনী।
\end{proof}

\begin{prob}
  $\Mult$-এর আদিম পুনরাবৃত্ত সংজ্ঞাটিকে
  \olref[cmp][rec][prf]{defn:primitive-recursion}-এর আবশ্যক বিন্যাসে
  লেখা যায় এবং সংশ্লিষ্ট অপেক্ষক $f$ ও~$g$ আদিম পুনরাবৃত্ত—এগুলি
  দেখিয়ে \olref[cmp][rec][prf]{prop:mult-pr} প্রমাণ করুন।
\end{prob}

\begin{ex}
আদিম পুনরাবৃত্ত সংজ্ঞার আমাদের একেবারে প্রথম উদাহরণটি ছিল:
\begin{align*}
h(0) & =  1 \\
h(y+1) & =  2 \cdot h(y).
\end{align*}
এই অপেক্ষকটি \olref{defn:primitive-recursion}-এর আবশ্যক বিন্যাসে বসে না,
কারণ $k=0$। সংজ্ঞাটিতে ধ্রুবক $1$ ও $2$-ও আছে। প্রথম সমস্যাটি এড়াতে
একটি অকেজো আর্গুমেন্ট যোগ করে অপেক্ষক~$h'$ সংজ্ঞায়িত করা যাক:
\begin{align*}
h'(x_0, 0) & =  f(x_0) = 1 \\
h'(x_0, y+1) & =  g(x_0, y, h'(x_0, y)) = 2 \cdot h'(x_0, y).
\end{align*}
অপেক্ষক $f(x_0) = 1$-কে মিশ্রণের সাহায্যে $\Succ$ ও $\Zero$ থেকে
সংজ্ঞায়িত করা যায়: $f(x_0) = \Succ(\Zero(x_0))$। অপেক্ষক $g$-কে
$g'(z) = 2 \cdot z$ ও অভিক্ষেপগুলির মিশ্রণ দিয়ে সংজ্ঞায়িত করা যায়:
\begin{align*}
g(x_0, y, z) & = g'(\Proj{3}{2}(x_0, y, z))
\intertext{এবং $g'$-কে আবার মিশ্রণের সাহায্যে এভাবে সংজ্ঞায়িত করা যায়:}
g'(z) & = \Mult(g''(z), \Proj{1}{0}(z))
\intertext{এবং}
g''(z) & = \Succ(f(z)),
\end{align*}
যেখানে $f$ আগেরটিই: $f(z) = \Succ(\Zero(z))$। এখন~$h'$ পাওয়ার পর
আবার মিশ্রণ ব্যবহার করে $h(y) =
h'(\Proj{1}{0}(y),\Proj{1}{0}(y))$ নেওয়া যায়। এতে দেখা যায়, মৌলিক
অপেক্ষকগুলি থেকে ধারাবাহিক কয়েকটি মিশ্রণ ও আদিম পুনরাবৃত্তি ব্যবহার
করে $h$-কে সংজ্ঞায়িত করা যায়; অতএব $h$~আদিম পুনরাবৃত্ত।
\end{ex}

\end{document}
